\section{Разработка ПО для постобработки данных}

Для получения геометрических параметров железнодорожного пути необходимо выполнить обработку записанных массивов инерциальных данных с применением навигационного алгоритма, описанного в разделе \ref{sec:nav_algoritm}. С этой целью в рамках проекта было разработано специализированное программное обеспечение для постобработки данных.

Программный комплекс реализует следующие функции:

\begin{itemize}
	\item загрузка и анализ бинарных логов, сформированных микроконтроллером во время измерительного проезда;
	\item выполнение полного навигационного расчёта траектории в соответствии с математическим аппаратом, изложенным в разд. 5.7;
	\item визуализацию результатов в виде графиков временных рядов полученных и рассчитанных массивов величин;
	\item формирование структурированных отчётов для последующего анализа и паспортизации участков пути.
\end{itemize}

Архитектура приложения построена по модульному принципу, что позволяет легко добавлять новые методы обработки и адаптировать алгоритмы под требуемые задачи. Реализация выполнена на языке Python с применением готовых библиотек (NumPy, Matplotlib), что обеспечивает гибкость разработки и удобство сопровождения кода.